\documentclass{article}
\usepackage[left=2cm, right=2cm, top=2cm, bottom=2cm]{geometry}
\usepackage{graphicx}
\usepackage{amsmath}
\usepackage{amssymb}
\usepackage{amsthm}
\usepackage{fancyhdr}
\usepackage{verbatim}
\usepackage{listings}
\usepackage{xcolor}
\usepackage{pdfpages}
\usepackage{cancel}
\usepackage{physics}
\usepackage{esint}

\lstset{
    language=C++,
    basicstyle=\ttfamily\footnotesize,
    keywordstyle=\color{blue},
    commentstyle=\color{green},
    stringstyle=\color{red},
    numbers=left,
    numberstyle=\tiny\color{gray},
    frame=single,
    breaklines=true,
    captionpos=b,
}


\title{MTH 553: Lecture \# 26}
\author{Cliff Sun}

\newtheorem{theorem}{Theorem}[section]
\newtheorem{lemma}[theorem]{Lemma}
\newtheorem{definition}[theorem]{Definition}
\newtheorem{conjecture}[theorem]{Conjecture}
\newtheorem{proposition}[theorem]{Proposition}
\newtheorem{corollary}[theorem]{Corollary}
\newtheorem{one minute paper}[theorem]{One Minute Paper}

\pagestyle{fancy}
\lhead{\textbf{\thepage}\ \ \nouppercase{\rightmark}}
\chead{MTH 553: Lecture \# 26}
\rhead{Cliff Sun}

\begin{document}

\maketitle

\section*{Lecture Span}
\begin{enumerate}
    \item Green's functions? Missed first 10 mins of class
\end{enumerate}

The green's functions

\begin{theorem}
    (Green 1) 
    \begin{equation}
        \int_{\Omega}(\nu \nabla u + \nabla \nu \nabla u) dx = \int_{\partial\Omega}v \frac{\partial u}{\partial \nu}dS
    \end{equation}
\end{theorem}

\begin{theorem}
    (Green 2)
    \begin{equation}
        \int_{\Omega}(\nu \triangle u - u \triangle \nu) dx = \int_{\partial \Omega}\left(\nu \partial_{\nu}u - u \partial_{u}\nu \right)
    \end{equation}
\end{theorem}

Both can be proved using the Divergence theorem.

Corollary, 

\begin{equation}
    \int_{\Omega}\triangle u dx = \int_{\partial\Omega}\frac{\partial u}{\partial \nu}dS
\end{equation}

By taking $v \equiv 1$. 

\begin{theorem}
    The Dirichlet problem for Poisson's equation is 
    \begin{align}
        -\triangle u &= f, \quad \text{in $\Omega$} \\
        u &= g, \quad \text{ on $\partial \Omega$}
    \end{align}
    Has at most one solution
\end{theorem}

\begin{theorem}
    The Neumann problem for Poisson's equation is 
    \begin{align}
        -\triangle u &= f, \quad \text{in $\Omega$} \\
        u &= g, \quad \text{ on $\partial \Omega$}
    \end{align}
    Has a solution only if 
    \begin{equation}
        -\int_\Omega f dx = \int_{\partial \Omega}h dx
    \end{equation}
    Here, this is the same as saying the heat source inside the region is balanced by the outwards heat flux at the boundary. There is one solution $u \in C^{2}(\overline{\Omega})$ up 
    to an additive constant. In other words, $u+c$ is also a solution if $c$ is a constant. 
\end{theorem}

\begin{proof}
    Part (a), suppose $u_1$ and $u_2$ are solutions. Then, 
    \begin{align}
        \triangle (u_1 - u_2) &= 0 \\
        u_1 - u_2 &= 0, \quad \text{ on $\partial \Omega$}
    \end{align}
    Then, applying the first greens theorem, then we obtain 
    \begin{equation}
        \int_{\Omega}\nabla |(u_1 - u_2)|^2 dx = 0
    \end{equation}
    With $u = u_1 - u_2$ and $v = u_1 - u_2$. Then 
    \begin{equation}
        \nabla (u_1 - u_2) \equiv 0 \implies u_1 - u_2 = \text{const.}
    \end{equation}
    Using the fact that $\Omega$ is connected. 
\end{proof}

Remark, the Cauchy problem for Poisson's equation is ill posed on $\Omega$. This is where $u$ and $\partial_{\nu}u$ are specified. This is two degrees of freedom but uniqueness implies only one degree of freedom. 

\subsection*{Mean Values and Maximum Principle}

New notation: 

\begin{equation}
    \fint_{\partial B(x,r)} u dS = \frac{1}{|\partial B(x,r)|}\int_{\partial B(x,t)}u(y)dy
\end{equation}
\begin{equation}
    = \frac{1}{\omega_n}\int_{S^{n-1}}u(x + r\xi)d\xi
\end{equation}

This is the mean value of $u$ over the surface of the ball $B(x,r)$, centered at $x$ with radius $r$. 

Moreover, 

\begin{equation}
    \fint_{B(x,r)}udy 
\end{equation}

is the mean value of $u$ over the ball. 

\begin{theorem}
    Mean Value Theorem. Let $u \in C^{2}(\Omega)$, then $u$ is harmonic 
    \begin{equation}
        \iff u(x) = \fint_{\partial B(x,r)}udS
    \end{equation}
    For all $x,t$ with $B(x,r) \subset \Omega$. Similarly, 
    \begin{equation}
        u(x) = \fint_{B(x,r)}udy
    \end{equation}
\end{theorem}

\begin{proof}
    Proof of MVT (mean value theorem). Calculate
    \begin{equation}
        \int_{B(x,r)}\triangle u dy = \int_{\partial B(x,r)}\frac{\partial u}{\partial \nu}dS
    \end{equation}
    by the corollary. 
    \begin{align*}
        &\int_{S^{n-1}}\left[ \partial_r u(x + r\xi) \right]r^{n-1}d\xi \\
        &= \omega_n r^{n-1}\partial_r\left[ \frac{1}{\omega_n} \int_{S^{n-1}}u(x + r\xi)d\xi\right] \\
        &= \omega_n r^{n-1}\partial_r \fint_{\partial B(x,r)}u dy
    \end{align*}
    Here, $\nu$ is the normal derivative and which is the same as $r$. 
    This is a very important calculation. If $\triangle u =0 $, then this mean value doesn't change. And in particular, 
    \begin{equation}
        \fint_{\partial B(x,r)}u dy = \text{const.}
    \end{equation}
    For all $r$. Take $r$ to be really small, then this is simply just $u(x)$. 
\end{proof}

Conversely, if $\triangle u \neq 0$. Then 
\begin{equation}
    \partial_{r} \fint_{\partial B(x,r)}udS \neq 0
\end{equation}
And in particular, $\partial_r [\dots] > 0$ if $\triangle u > 0$ on some ball $B(x,\rho)$. Hence, you don't get equality anymore. Therefore, $u$ is harmonic if it satisfies the mean value theorem. 
\end{document}